\subsection{Objetivo e arquitetura da solução}

O componente de inteligência artificial atua como um engenheiro de pista virtual, interpretando a corrida para produzir recomendações em português, indicar sua urgência e sugerir ações como atacar, defender, conservar equipamento ou preparar um pit stop. Seu trabalho não se limita à redação das mensagens: é necessário identificar o piloto acompanhado, selecionar acontecimentos relevantes e apresentar a orientação enquanto ela ainda fizer sentido. Para o jogador, a IA tem caráter consultivo; a atuação automática é reservada aos competidores controlados pela simulação.

Três fluxos devem ser distinguidos: \texttt{agent.py} analisa voltas de corridas históricas; \texttt{redis\_listener.py}, iniciado pelo container \texttt{ai}, processa eventos ao vivo com modelo de linguagem; e \texttt{ai\_engineer.py}, no backend, aplica regras autônomas e pode interferir nos demais carros. Eles compartilham eventos e formatos de recomendação, mas a análise histórica não é acionada pelo listener, e o engenheiro do backend independe do Ollama.

No fluxo ao vivo, o Redis aplica Publish/Subscribe para desacoplar simulador, IA e interfaces. O backend publica em \texttt{race:\{id\}:events}; o listener mantém um estado por corrida e devolve sugestões ao canal, encaminhadas pela FastAPI via WebSocket. O piloto acompanhado é obtido dos metadados, com prioridade para \texttt{followed\_driver\_id} e \texttt{user\_driver\_id}, antes do valor padrão de ambiente. Assim, a seleção da interface orienta a análise, e o \texttt{race\_id} separa corridas simultâneas. Modelo, IA, backend e visualizadores também são processos independentes no Compose.

\subsection{Contexto e inferência com Ollama}

A classe \texttt{LiveRaceState} funciona como um objeto de contexto em memória, reunindo volta, clima, posições e avisos recentes de pneus e pit stops. Cada evento atualiza essa estrutura e pode revelar mudanças de posição, alterações climáticas ou saídas e retornos de pista. Quando há pelo menos quatro carros, incidentes de outros pilotos são filtrados pela proximidade com o competidor acompanhado, evitando que um acontecimento distante tenha o mesmo peso de um risco imediato. Além desses gatilhos, atualizações periódicas em tempo de simulação permitem reavaliar a estratégia durante trechos estáveis.

O prompt resume essa situação e exige JSON com aviso, previsão de pit stop e urgência. A validação confere os campos obrigatórios, completa previsões incompletas e normaliza urgências desconhecidas. O módulo \texttt{client.py} atua como adaptador HTTP para o endpoint \texttt{/api/generate}, isolando os detalhes do Ollama. Modelo e timeout podem ser configurados por ambiente, com \texttt{llama3.2} e 30 segundos como padrão; \texttt{keep\_alive}, contexto reduzido e limite de geração buscam diminuir a latência e a variabilidade. Ainda assim, a chamada é síncrona e pode atrasar o processamento do evento seguinte.

Na análise histórica, o pacote \texttt{context} consulta \texttt{data/processed/f1.db} para reunir classificação, diferenças de tempo, voltas recentes, pneus, clima e safety cars. As estimativas de desgaste são indícios, não previsões determinísticas. Esse caminho apoia análises e testes sobre dados conhecidos, mas não implica que essas consultas estejam integradas ao listener ao vivo.

\subsection{Decisão tática e tolerância a falhas}

O engenheiro do backend examina periodicamente o \textit{snapshot} do simulador e relaciona fase da corrida, janela de pit, pneus, combustível e diferenças entre pilotos. Rival próximo à frente favorece \texttt{overtake} ou \texttt{push}; pressão recebida leva a \texttt{defend}; desgaste e combustível insuficiente motivam \texttt{conserve}; e liderança confortável permite administrar o ritmo. Cada recomendação registra motivo, urgência e confiança, tornando a decisão mais explicável do que uma simples mudança de comportamento do carro.

A execução é separada da análise pelo \texttt{CommandBus}, que guarda no Redis o comando recente de cada piloto e também atende ao endpoint manual de comandos da API. O \texttt{CoachController} envolve o controlador original segundo o padrão Decorator, consome comandos ainda não processados e ajusta entradas como aceleração e frenagem. Eventos de comando e sugestão expõem essa atuação às interfaces. No carro do jogador, porém, o engenheiro publica somente a recomendação textual, sem enviar a ordem ao controle.

Quando o Ollama falha ou responde com JSON inválido, \texttt{fallback\_tactical.py} seleciona uma resposta determinística para situações como últimas voltas, safety car, rivais próximos, pneus críticos e pista molhada. O resultado conserva o contrato da sugestão e é identificado como fallback. Um mecanismo de antirrepetição complementa essa proteção: eventos reais permitem uma nova fala quando comando ou texto muda, enquanto gatilhos periódicos silenciam repetições. Como a comparação é mantida por corrida, o mecanismo reduz mensagens redundantes sem misturar o estado de simulações diferentes.

\subsection{Persistência, apresentação e observabilidade}

SQLite, Redis e memória têm papéis distintos. O primeiro conserva o histórico consultado pelos contextos; o segundo armazena metadados, snapshots e comandos operacionais, além de transportar eventos; e a última mantém o estado recente do listener e da antirrepetição. Esse estado transitório não é restaurado após reinicialização, e a publicação em um canal não equivale a um histórico persistente com replay garantido.

No frontend, parser e store convertem eventos em sugestões, comandos e indicadores de ciclo de vida, validando o \texttt{raceId} quando disponível. O painel flutuante apresenta status, erros e recomendações nas visões 2D e 3D, complementado por uma representação tridimensional cuja cor e pulsação comunicam atividade e urgência. Resta um alinhamento: o parser exige \texttt{driver\_id} em \texttt{ai\_suggestion}, mas o listener do Ollama acrescenta somente tipo e volta. Como as sugestões do backend já identificam o piloto, é necessário padronizar os eventos do listener para evitar seu descarte pela interface.

A observabilidade foi organizada em \texttt{metrics.py}, com envio periódico previsto para o MLflow. As chamadas ao modelo alimentam medidas de latência, disponibilidade e erros, enquanto a validação e o listener registram JSON inválido e mensagens suprimidas. O coletor também prevê indicadores de fallback e percentis, permitindo avaliar confiabilidade e desempenho sem depender apenas da impressão causada pelas respostas.

\subsection{Testes e limitações}

Os testes em \texttt{ai/tests/} abrangem cliente Ollama, esquemas, contextos históricos, eventos ao vivo e antirrepetição; o backend possui testes para o controle do engenheiro e eventos de demonstração. Eles verificam regras isoladas, mas não substituem a validação integrada com Redis, modelo e visualizador, nem comprovam a qualidade estratégica de cada recomendação. Os resultados devem ser associados ao ambiente e à versão executados, pois parte dos testes históricos depende do conteúdo específico do banco SQLite.

Permanecem limitações relevantes: a espera síncrona pelo modelo pode atrasar eventos; o contexto ao vivo é perdido na reinicialização; não há recuperação completa após falha prolongada do Redis; e as regras autônomas privilegiam explicabilidade, sem garantir estratégia ótima. Também é necessário concluir o contrato entre listener e frontend. Por fim, Redis exposto sem autenticação explícita e inferência por HTTP local são adequados ao desenvolvimento, mas exigem restrição de acesso e controles adicionais antes de uma implantação pública.
